\documentclass[11pt]{article}
\usepackage[T1]{fontenc}
\usepackage[utf8]{inputenc}
\usepackage{lmodern,amsmath,amssymb,amsthm,mathtools}
\usepackage[margin=25mm]{geometry}
\usepackage[hidelinks]{hyperref}
\newcommand{\C}{\mathbb C}
\newcommand{\R}{\mathbb R}
\newcommand{\PP}{\mathcal P}
\newcommand{\norm}[1]{\left\lVert#1\right\rVert}
\DeclareMathOperator{\Tr}{Tr}
\DeclareMathOperator{\rank}{rank}
\DeclareMathOperator{\im}{im}
\DeclareMathOperator{\diag}{diag}
\DeclareMathOperator{\spec}{spec}
\newtheorem{theorem}{Theorem}
\newtheorem{proposition}[theorem]{Proposition}
\title{The full-primary positive-trace minimum and its original-metric receivers}
\author{Split-Zero mathematical continuation}
\date{19 September 2026}
\begin{document}
\maketitle

\begin{abstract}
For the unchanged arithmetic multiplication packet, the minimum of the
positive Hermitian weight trace over all positive metrics is evaluated,
attained, and its entire minimizing set is parametrized by finite
Lyapunov sums. All primary jets and all same-half cross terms are
retained. A rank-one observable forcing on each complete spectral half
constructs an attaining metric whose weight has precisely two nonzero
eigenvalues. The canonical attained polynomial metric is then evaluated
through its actual boundary rows: its exact rank-two scalar exceeds the
minimum by an explicit positive spectral-projection expression. The
positive-primary determinant line has exactly the same growth exponent
for every metric. Finally an original finite polynomial section realizes
every minimizing metric up to a specified scalar, and its complete
residual in each smaller polynomial source quotient is computed.
\end{abstract}

\section{Original packet, source and coordinate maps}

The incoming note \cite{incoming} proposes the minimum and its complete
Lyapunov parametrization. The present text proves them, adds the
rank-two attaining construction, and evaluates their precise receivers.
The fixed source and cutoff conventions are GMF1--3 of
\cite{gmf}; the earlier exact metric and source maps are TC1--29d of
\cite{tc}. These source references record provenance. Every new finite
assertion below is proved here.

Keep the original hypothetical actual quartet with full multiplicity:
\begin{equation}\tag{PT1}
\begin{gathered}
 \rho=\tfrac12+\delta+i\gamma,\quad
 0<\delta<\tfrac12,\quad\gamma>2,\quad m_0\ge1,\\
 h(s)=\prod_{\zeta\in\{\rho,\bar\rho,1-\rho,1-\bar\rho\}}
                     (s-\zeta)^{m_0},\quad
 v_h=2\xi/h,\quad w_h(y)=|v_h(1/2+iy)|^2/(2\pi),\\
 m_k=w_h^{*k},\quad k\equiv1\pmod4,\quad k\ge5,\quad
 e=1+k(m_0-1),\quad q=e(k+1)^2,\quad c=k/2,\\
 \lambda_{ab}=c+(2a-k)\delta+i(2b-k)\gamma,\quad0\le a,b\le k,\\
 \chi(S)=\prod_{a,b=0}^k(S-\lambda_{ab})^e,\quad
 E=\C[S]/(\chi),\quad A=M_S:E\to E,\quad B=A-cI.
\end{gathered}
\end{equation}
The original basis of \(E\) is \(1,S,\ldots,S^{q-1}\).
The numbers \(\lambda_{ab}\) are pairwise distinct: equality first
forces the real indices \(a\), and then the imaginary indices \(b\),
to coincide. Thus \(\chi\) has exactly the displayed complete primary
factors. None is divided out in this paper.

For \(G>0\) in this original basis define
\begin{equation}\tag{PT2}
 W(G)=A^*G+GA-kG=B^*G+GB,\quad
 H_G=G^{-1/2}W(G)G^{-1/2},\quad
 T_+(G)=\Tr(H_G)_+.
\end{equation}
The symbol \(X_+\) means the positive part of a Hermitian matrix,
obtained by replacing each negative eigenvalue by zero; write
\(X_-=(-X)_+\). Positive square roots always refer to the fixed
standard coordinate Hermitian form. The physical coordinate in the
source remains \(S=c+iy\); its measure is exactly \(m_k(y)\,dy\).

Here is the complete confluent coordinate map and its inverse. Define
\begin{equation}\tag{PT3}
 C:E\longrightarrow\bigoplus_{a,b=0}^k\C^e,\qquad
 (CP)_{(\lambda,r)}=P^{(r)}(\lambda)/r!,\qquad
 C_{(\lambda,r),j}=\begin{cases}\binom jr\lambda^{j-r},&j\ge r,\\0,&j<r.
 \end{cases}
\end{equation}
For each \(\lambda\) put \(h_\lambda=\chi/(S-\lambda)^e\) and
write the actual Taylor series
\(h_\lambda(\lambda+z)^{-1}=\sum_{n\ge0}a_{\lambda,n}z^n\).
Its constant denominator \(h_\lambda(\lambda)\) is the full product
of the other root differences, each with exponent \(e\), and is
nonzero. Set
\begin{equation}\tag{PT4}
 p_{\lambda,r}(S)=h_\lambda(S)(S-\lambda)^r
       \sum_{n=0}^{e-r-1}a_{\lambda,n}(S-\lambda)^n,
 \qquad0\le r<e.
\end{equation}
This polynomial has degree at most \(q-1\). At \(\lambda\) its
Taylor class modulo \((S-\lambda)^e\) is \((S-\lambda)^r\);
at every other root its first \(e\) Taylor coefficients vanish.
Consequently these polynomials are the columns of \(C^{-1}\).
This proves both inverse laws of (PT3), without replacing any primary
factor by a simple root. Multiplication by \(S\) now gives
\begin{equation}\tag{PT5}
 CAC^{-1}=\bigoplus_\lambda(\lambda I_e+J_e),\qquad
 (J_e)_{r+1,r}=1\quad(0\le r<e-1),\quad J_e^e=0.
\end{equation}
The direct sum and the subsequent maps are coordinate morphisms;
the recorded operator remains the original \(A\).

Let \(E_+\) be the sum of the complete primary spaces with
\(a\ge(k+1)/2\), and \(E_-\) its complementary sum. Both have
dimension \(r=q/2\). Formulae (PT3)--(PT4) give explicit column
frames for their inclusions. For a column frame \(X_+:\C^r\to E_+\)
define the actual Euclidean orthogonal projection
\begin{equation}\tag{PT6}
 P_G=G^{1/2}X_+(X_+^*GX_+)^{-1}X_+^*G^{1/2}
       :\C^q\longrightarrow G^{1/2}E_+.
\end{equation}
It is self-adjoint and idempotent by matrix multiplication, and its
image is precisely the stated subspace. It is independent of the
choice of frame within \(E_+\).

\section{The trace bound, exact defect and equality}

The finite projection argument below is the positive-part form of
Ky Fan's variational trace principle \cite{fan}. We give the full
argument rather than use it as a black box.

\begin{theorem}
For every original-coordinate \(G>0\),
\begin{equation}\tag{PT7}
 L:=\frac{\delta q(k+1)}2=\Tr(P_GH_G),\qquad
 T_+(G)-L
 =\Tr((I-P_G)(H_G)_+)+\Tr(P_G(H_G)_-)\ge0.
\end{equation}
The two terms in the defect are exactly
\begin{equation}\tag{PT8}
 T_+(G)-L
 =\norm{(H_G)_+^{1/2}(I-P_G)}_F^2
   +\norm{(H_G)_-^{1/2}P_G}_F^2.
\end{equation}
Equality holds precisely when \(E_+\) and \(E_-\) are
\(G\)-orthogonal and their actual restricted metrics satisfy
\begin{equation}\tag{PT9}
 B_+^*G_++G_+B_+\succeq0,\qquad
 C_-^*G_-+G_-C_-\succeq0,
 \quad B_+=A|_{E_+}-cI,\quad C_-=cI-A|_{E_-}.
\end{equation}
\end{theorem}
\begin{proof}
Put \(K=G^{1/2}BG^{-1/2}\), so \(H_G=K+K^*\).
The subspace \(G^{1/2}E_+\) is \(K\)-invariant, and the
compression of \(K\) to it is similar to \(B_+\). Therefore
\[
 \Tr(P_GH_G)=2\Re\Tr B_+
 =e(k+1)2\delta\sum_{a=(k+1)/2}^k(2a-k).
\]
The summands in the last sum are the odd integers
\(1,3,\ldots,k\); their sum is \(((k+1)/2)^2\). This proves
the first equality in (PT7). Substitute
\(H_G=(H_G)_+-(H_G)_-\) and use linearity of trace to obtain
the second equality. For a positive matrix \(X\) and an orthogonal
projection \(P\),
\(\Tr(PX)=\Tr(PX^{1/2}X^{1/2}P)=\norm{X^{1/2}P}_F^2\).
Apply this twice to get (PT8) and nonnegativity, without assuming
that either projection commutes with either positive part.

The defect vanishes exactly when
\((H_G)_+^{1/2}(I-P_G)=0\) and \((H_G)_-^{1/2}P_G=0\).
These equalities say that every positive eigenvector of \(H_G\)
belongs to \(\im P_G\), and every negative eigenvector belongs
to its orthogonal complement. The zero eigenspace may be split in
any way. In particular \(P_GH_G=H_GP_G\).
In an orthonormal frame adapted to \(\im P_G\), invariance gives
\[
 K=\begin{pmatrix}K_+&D\\0&K_-\end{pmatrix},\qquad
 H_G=\begin{pmatrix}K_++K_+^*&D\\D^*&K_-+K_-^*\end{pmatrix}.
\]
Commutation forces \(D=0\). Thus the orthogonal complement is
also \(K\)-invariant; its spectrum is precisely the negative-real
part of the original spectrum. To identify this complement exactly,
apply the complete spectral idempotents supplied by (PT4) with
\(r=0\). A vector in this invariant complement is killed by the
product of the negative primary factors, so all positive idempotent
components vanish. Dimension \(r\) then proves that the complement
is \(G^{1/2}E_-\). The restrictions of \(H_G\) on the two halves
are respectively nonnegative and nonpositive, which is exactly
(PT9) by congruence with their positive square roots.
Conversely, orthogonality and (PT9) give those two signs on the
two invariant halves. The positive trace is then their positive
half trace \(2\Re\Tr B_+=L\). This proves both directions,
including all possible zero eigenvalues within either half.
\end{proof}

\section{An explicit full-primary attainer}

Put \(t_0=\delta/2\), \(D_e=\diag(1,t_0,\ldots,t_0^{e-1})\),
and define in the original frame
\begin{equation}\tag{PT10}
 T=\left(\bigoplus_\lambda D_e\right)C,\qquad G_*=T^*T.
\end{equation}
The map \(T:(E,G_*)\to(\C^q,I)\) is an exact isometry.
Indeed \(U=TG_*^{-1/2}\) is unitary, and directly
\[
 UH_{G_*}U^*=TAT^{-1}+(TAT^{-1})^*-kI.
\]
Since \(D_eJ_eD_e^{-1}=t_0J_e\), each block on the right is
\(2(2a-k)\delta I_e+t_0(J_e+J_e^*)\). The vector with
\(r\)-th entry \(\sin((r+1)j\pi/(e+1))\), \(0\le r<e\),
is an eigenvector of \(J_e+J_e^*\) with eigenvalue
\(2\cos(j\pi/(e+1))\): the first and last missing entries are
zero, and the interior recurrence is the sine addition identity.
These \(e\) eigenvalues are distinct, giving a full eigenbasis.
Consequently
\begin{equation}\tag{PT11}
 \spec H_{G_*}
 =\left\{2(2a-k)\delta+\delta\cos\frac{j\pi}{e+1}:
 0\le a,b\le k,\ 1\le j\le e\right\}.
\end{equation}
Every one of these values has the sign of \(2a-k\), since
\(|2a-k|\ge1\) and \(|\cos|<1\). The cosine terms sum to
zero in each block by pairing \(j\) with \(e+1-j\). Thus
\begin{equation}\tag{PT12}
 \min_{G>0}T_+(G)=T_+(G_*)=L,\quad
 \rank(H_{G_*})_+=\rank(H_{G_*})_-=q/2,\quad
 \Tr H_{G_*}=0,\quad\Tr|H_{G_*}|=2L.
\end{equation}
In particular \(\norm{H_{G_*}}=2k\delta\) when \(e=1\),
and \(2k\delta+\delta\cos(\pi/(e+1))\) when \(e>1\).
These are norms of this displayed metric. They do not replace
the positive trace of any other original metric.

\section{Every minimizing metric and every nilpotent term}

The method here is the classical Lyapunov matrix equation method,
named for A.~M.~Lyapunov \cite{lyapunov}. All convergence,
positivity and inverse statements needed here are proved below.
Use the fixed primary bases (PT3) on each half and put
\begin{equation}\tag{PT13}
 D_+=\bigoplus_{a>(k/2),b}( (\lambda_{ab}-c)I_e+J_e),\qquad
 D_-=\bigoplus_{a<(k/2),b}( (c-\lambda_{ab})I_e-J_e).
\end{equation}
Every eigenvalue of \(D_+\) and \(D_-\) has strictly positive
real part. Notice the minus sign on every negative-half nilpotent.

\begin{theorem}
All minimizing metrics, with no duplicates suppressed, are obtained
by choosing \(Q_+,Q_-\succeq0\) such that
\begin{equation}\tag{PT14}
 \bigcap_{j=0}^{r-1}\ker(Q_\pm^{1/2}D_\pm^j)=\{0\},
 \qquad r=q/2,
\end{equation}
and setting
\begin{equation}\tag{PT15}
 \Gamma_\pm=\int_0^\infty e^{-tD_\pm^*}Q_\pm e^{-tD_\pm}\,dt,
 \qquad G=C^*\diag(\Gamma_+,\Gamma_-)C.
\end{equation}
In this formula \(C\) is ordered first by the positive half and
then by the negative half, with the same primary coordinates
within each half. Every integral is a finite rational block sum.
For blocks \(D_\alpha=b_\alpha I+N_\alpha\) its exact value is
\begin{equation}\tag{PT16}
 \Gamma_{\alpha\beta}
 =\sum_{i,j=0}^{e-1}
 \frac{(-1)^{i+j}(i+j)!}
      {i!j!(\bar b_\alpha+b_\beta)^{i+j+1}}
 (N_\alpha^*)^i Q_{\alpha\beta}N_\beta^j.
\end{equation}
All blocks within each half, including off-diagonal blocks, occur.
\end{theorem}
\begin{proof}
The full block exponential is the finite sum
\(e^{-tb_\alpha}\sum_{j=0}^{e-1}(-t)^jN_\alpha^j/j!\).
Every entry of the integrand in (PT15), and its derivative, is
therefore a polynomial times an exponential of strictly negative
real exponent. The integrals converge absolutely and the endpoint
at infinity is zero. Differentiation and integration give
\[
 D_\pm^*\Gamma_\pm+\Gamma_\pm D_\pm=Q_\pm.
\]
The quadratic form of \(\Gamma_\pm\) on \(v\) is the integral
of \(\norm{Q_\pm^{1/2}e^{-tD_\pm}v}^2\). It vanishes
exactly when this continuous nonnegative integrand is identically
zero. Its entire power series is zero exactly when all
\(Q_\pm^{1/2}D_\pm^jv\) vanish. The characteristic polynomial
of \(D_\pm\) has degree \(r\); division by it expresses every
power \(j\ge r\) as a linear combination of lower powers.
This proves the finite test (PT14) and strict positivity.

Conversely let \(\Gamma>0\) on one half satisfy
\(Q=D^*\Gamma+\Gamma D\succeq0\). Differentiating
\(e^{-tD^*}\Gamma e^{-tD}\), and integrating between zero
and infinity, recovers (PT15) exactly. The preceding positivity
argument proves (PT14). The same argument shows that the
homogeneous Lyapunov equation has only the zero solution.
The characterization (PT9) now proves that (PT15) gives every
minimizer and only minimizers.

For (PT16), expand both finite exponentials. The term with powers
\(i,j\) has scalar integral
\(\int_0^\infty t^{i+j}e^{-zt}\,dt=(i+j)!/z^{i+j+1}\),
where \(z=\bar b_\alpha+b_\beta\) has positive real part.
The integral identity starts with \(1/z\) and follows by
successive integration by parts, whose endpoints vanish by the
same exponential decay. This proves every coefficient and sign.
On the negative half \(N_\alpha=-J_e\), so its two powers
supply an additional \((-1)^{i+j}\); these signs remain
explicit in (PT13)--(PT16). No inverse of a nilpotent is used.
\end{proof}

\section{A minimizing metric with rank-two weight}

There is an attaining family with many zero weight directions,
even when \(e>1\). On each half define the single literal row
\begin{equation}\tag{PT17}
 \ell_\pm(v)=\sum_{\lambda\text{ in that half}}v_{\lambda,e-1},
 \qquad Q_\pm=\ell_\pm^*\ell_\pm.
\end{equation}
The row extracts the final divided-derivative coefficient in every
complete primary block, with coefficient exactly one in each.

\begin{proposition}
The two choices (PT17) satisfy (PT14). The metric
\(G_{\mathrm{two}}\) obtained from (PT15)--(PT16) has
\begin{equation}\tag{PT18}
 \spec H_{G_{\mathrm{two}}}=\{L,-L,0,\ldots,0\},\qquad
 T_+(G_{\mathrm{two}})=\norm{H_{G_{\mathrm{two}}}}=L.
\end{equation}
Hence the positive-trace minimum is attained within the rank-two
weight stratum as well as within the full-rank weight stratum.
\end{proposition}
\begin{proof}
Write the blocks of either half as \(b_\lambda I+sJ_e\), with
the fixed common sign \(s=1\) on the positive half and \(s=-1\)
on the negative half. Then
\[
 \ell e^{-tD}v
 =\sum_\lambda e^{-tb_\lambda}
      \sum_{j=0}^{e-1}\frac{(-st)^j}{j!}v_{\lambda,e-1-j}.
\]
Distinct \(\lambda\)'s give distinct \(b_\lambda\)'s on
either half. The functions \(e^{-tb_\lambda}p_\lambda(t)\),
with \(\deg p_\lambda<e\), are linearly independent. Here is
an explicit proof: apply
\(\prod_{\mu\ne\lambda}(\partial_t+b_\mu)^e\) to a
vanishing sum. It kills every summand except the \(\lambda\)
summand, where it gives \(e^{-tb_\lambda}\) times
\(\prod_{\mu\ne\lambda}(\partial_t+b_\mu-b_\lambda)^e
p_\lambda\). On the polynomials of degree less than \(e\),
each operator \(\partial_t+a\) with \(a\ne0\) is invertible:
its matrix is triangular with diagonal \(a\), or its inverse
is the finite geometric sum in \(-\partial_t/a\).
Therefore every \(p_\lambda\) is zero. Their displayed
coefficients force every coordinate of \(v\) to be zero.
Thus \(\ell e^{-tD}v=0\) for all \(t\) implies \(v=0\),
and (PT14) follows as in the preceding theorem.

In each half the Lyapunov weight is \(Q_\pm\), of rank one.
Whitening by the positive \(\Gamma_\pm\) preserves this rank
and gives a nonnegative rank-one matrix on the positive half
and a nonpositive rank-one matrix on the negative half. Its
nonzero eigenvalue is its trace. Cyclicity of trace gives
\[
 \Tr(\Gamma_+^{-1}Q_+)=2\Re\Tr D_+=L,
 \qquad\Tr(\Gamma_-^{-1}Q_-)=2\Re\Tr D_-=L.
\]
The two complete halves are orthogonal in the constructed metric,
so these are exactly the eigenvalues in (PT18). All vectors of
all primary blocks remain present; zero Hermitian weight does not
remove those vectors or their original operator action.
\end{proof}

The rank-two metric has fully evaluated entries as well. On either
half put \(\sigma=1\) for \(D_+\), \(\sigma=-1\) for \(D_-\),
and, for indices \(0\le i,j<e\), put \(u=e-1-i\), \(v=e-1-j\).
In (PT16), \(Q_{\alpha\beta}\) is the matrix with its only nonzero
entry, equal to one, at \((e-1,e-1)\). Thus
\begin{equation}\tag{PT18a}
 (\Gamma_\pm)_{\alpha i,\beta j}
 =\frac{(-\sigma)^{u+v}(u+v)!}
        {u!v!(\bar b_\alpha+b_\beta)^{u+v+1}}.
\end{equation}
For \(u+v>0\), substituting these entries into the Lyapunov
equation reduces its numerator to
\[
 \frac{(u+v)!}{u!v!}
 -\mathbf1_{u>0}\frac{(u+v-1)!}{(u-1)!v!}
 -\mathbf1_{v>0}\frac{(u+v-1)!}{u!(v-1)!}=0.
\]
For \(u=v=0\) the single term gives one. This directly verifies
every row and both nilpotent signs independently of the integral.

Every positive/negative weight-rank pair
\(1\le r_+,r_-\le q/2\) occurs among minimizers. In each half
choose a positive semidefinite matrix of the prescribed rank whose
range contains \(\ell_\pm^*\). Its kernel is contained in
\(\ker\ell_\pm\), so its finite observability kernel is zero
by the proof of (PT17). Formula (PT15) gives a positive metric,
and invertible congruence preserves the prescribed weight rank.
This proves the claim, including all intermediate zero-weight
dimensions, with every original primary vector present.

\section{The original positive-primary exterior line}

Equip every exterior power with the Gram-determinant metric induced
by its stated \(G\). The inclusion of the positive-primary line is
the actual map
\(\Lambda^rX_+:\Lambda^r\C^r\to\Lambda^r E\), not a
choice of a positive eigenspace of \(H_G\). On this line the
generator induced by \(B\) is multiplication by \(\Tr B_+\).
Its imaginary part vanishes because for each fixed \(a\) the
sum of \(2b-k\) over \(0\le b\le k\) is zero. Its real
part is \(L/2\) by (PT7). Thus for any nonzero element
\(\omega_+\in\Lambda^r E_+\) and every real \(t\),
\begin{equation}\tag{PT19}
 \norm{\Lambda^r e^{tB}\omega_+}_{\Lambda^rG}^{,2}
       =e^{Lt}\norm{\omega_+}_{\Lambda^rG}^{,2},\qquad
 \norm{\Lambda^r e^{tA}\omega_+}_{\Lambda^rG}^{,2}
       =e^{(kr+L)t}\norm{\omega_+}_{\Lambda^rG}^{,2}.
\end{equation}
To verify the claimed action directly, triangularize each full
primary block by its fixed basis (PT5). The determinant of
\(e^{tB_+}\) is the product of its diagonal entries,
\(e^{t\Tr B_+}\); all nilpotent exponential diagonals are one.
Squaring the scalar modulus gives (PT19). In the original column
frame the derivative has the equally exact expression
\begin{equation}\tag{PT20}
 \left.\frac{d}{dt}\right|_{t=0}
 \log\det\bigl(X_+^*e^{tB^*}Ge^{tB}X_+\bigr)
 =\Tr\bigl((X_+^*GX_+)^{-1}X_+^*W(G)X_+\bigr)
 =\Tr(P_GH_G)=L.
\end{equation}
The derivative of \(\log\det M(t)\) follows by multilinear
expansion of \(\det(I+tM(0)^{-1}M'(0))\). Substituting (PT6)
and moving finite factors around the trace proves the second
equality. This supplies the exact receiver to the original exterior
benchmark for every \(G\), including the fixed attained metric.

\section{The canonical attained metric and its exact rank-two gap}

Use the original source \(m_k(y)\,dy\) and physical coordinate
\(S=c+iy\) from (PT1). Its positivity and moments are the original
source data of \cite{gmf}; in particular it is positive on a set
of positive measure and all polynomial moments are finite. For
\(\PP_N=\C[S]_{\le N}\), \(N\ge q-1\), put
\begin{equation}\tag{PT21}
 \begin{gathered}
 \langle f,g\rangle=\int_\R\overline{f(c+iy)}g(c+iy)m_k(y)\,dy,
 \quad H_N=(\langle S^i,S^j\rangle)_{0\le i,j\le N},\\
 J_N:\PP_N\twoheadrightarrow E,\quad f\mapsto[f]_\chi,\qquad
 G_N=(J_NH_N^{-1}J_N^*)^{-1},\quad
 R_N=H_N^{-1}J_N^*G_N.
 \end{gathered}
\end{equation}
A nonzero polynomial has only finitely many zeros on the vertical
line; thus \(H_N>0\). Surjectivity of \(J_N\) gives \(G_N>0\).
Matrix multiplication gives \(J_NR_N=I\) and
\(R_N^*H_NR_N=G_N\). For \(z\in\ker J_N\),
\(z^*H_NR_N=z^*J_N^*G_N=0\). Therefore \(R_Nu\) is the
actual attained minimum representative of \(u\); every other
representative has squared norm \(u^*G_Nu+\norm z^2\).

Define the following original rows and polynomial, with no omitted
leading coefficient:
\begin{equation}\tag{PT22}
 \ell_N(u)=[S^N]R_Nu,\qquad
 r_N(S)=\chi(S)S^{N+1-q},\qquad
 s_N(u)=\langle r_N,R_Nu\rangle.
\end{equation}
The exponent \(N+1-q\) is nonnegative even at \(N=q-1\).
Set \(T_Nf=Sf-[S^N]f\,r_N\). As \(\chi\) is monic,
the degree-\(N+1\) coefficient cancels, so \(T_N:\PP_N\to\PP_N\).
It satisfies the exact quotient intertwining relation
\(J_NT_N=AJ_N\). Thus \(R_NA\) is the orthogonal projection
of \(T_NR_N\) onto \((\ker J_N)^\perp\).
The physical identity \(\bar S+S=k\) on the vertical line gives
\(\langle Sf,g\rangle+\langle f,Sg\rangle=k\langle f,g\rangle\).
Taking \(f=R_Nu,g=R_Nv\), and subtracting the two displayed
leading-coefficient corrections, proves
\begin{equation}\tag{PT23}
 W(G_N)=-(\ell_N^*s_N+s_N^*\ell_N),\qquad
 \rank W(G_N)\le2,\qquad
 \Tr H_{G_N}=2\Re\Tr A-kq=0.
\end{equation}
The last equality follows by summing \(2a-k\) over all \(a,b\).
Since (PT7) gives a strictly positive trace lower bound, the rank
is exactly two and its two nonzero eigenvalues are
\(+\epsilon_N,-\epsilon_N\), with \(\epsilon_N>0\).

Define, in the original packet coordinates,
\begin{equation}\tag{PT24}
 u_N=s_NG_N^{-1}s_N^*,\qquad
 v_N=\ell_NG_N^{-1}\ell_N^*,\qquad
 z_N=s_NG_N^{-1}\ell_N^*.
\end{equation}
The trace identity in (PT23) gives \(\Re z_N=0\). Put
\(a=G_N^{-1/2}\ell_N^*\) and \(b=G_N^{-1/2}s_N^*\),
so \(H_{G_N}=-(ab^*+ba^*)\). Direct multiplication yields
\[
 \Tr(H_{G_N}^2)=2u_Nv_N+2\Re(z_N^2)
                 =2u_Nv_N-2(\Im z_N)^2.
\]
The eigenvalues in (PT23) give the exact certificate and bound
\begin{equation}\tag{PT25}
 \epsilon_N^2=u_Nv_N-(\Im z_N)^2
             =u_Nv_N-|z_N|^2\ge L^2.
\end{equation}
No coefficient of \(R_N\), no root difference in \(\chi\),
and no term in either original source region was changed to obtain
this identity.

Choose orthonormal eigenvectors \(p_N,n_N\) of the whitened
canonical matrix with eigenvalues \(\epsilon_N,-\epsilon_N\).
Then the full exact receiver of (PT7) is
\begin{equation}\tag{PT26}
 \epsilon_N-L
 =\epsilon_N\left(1-\norm{P_{G_N}p_N}^2
                         +\norm{P_{G_N}n_N}^2\right).
\end{equation}
This follows by substituting
\((H_{G_N})_+=\epsilon_Np_Np_N^*\) and
\((H_{G_N})_-=\epsilon_Nn_Nn_N^*\) in (PT7).
In particular equality means precisely that \(p_N\) lies in
\(G_N^{1/2}E_+\) and \(n_N\) lies in its orthogonal complement;
equivalently (PT9) holds for the fixed \(G_N\).
Equations (PT22)--(PT26) evaluate the comparison through its actual
minimum map. They do not assert that this minimum sequence happens
to lie in the minimizing set (PT15).

\section{Exact metric realization in the original finite source}

Use \(N=2q-1\), the unmodified \(H_N,J_N,R_N,G_N\) in
(PT21), and the complete relation map
\begin{equation}\tag{PT27}
 K_\chi:\PP_{q-1}\hookrightarrow\PP_{2q-1},\quad
 f\mapsto\chi f,\qquad
 Q_\chi=K_\chi^*H_NK_\chi>0,\quad
 B_\chi=K_\chi Q_\chi^{-1/2}.
\end{equation}
Injectivity follows from the integral domain \(\C[S]\).
Polynomial division gives \(\im K_\chi=\ker J_N\), exactly,
and this space has dimension \(q\). Hence
\(J_NB_\chi=0\), \(B_\chi^*H_NB_\chi=I\), and
\(B_\chi^*H_NR_N=0\).

For every displayed minimizing metric \(G_0\) from (PT15),
including \(G_*\) and \(G_{\mathrm{two}}\), set
\begin{equation}\tag{PT28}
 \begin{gathered}
 \alpha=1+\norm{G_0^{-1/2}G_NG_0^{-1/2}},\qquad
 \Delta=\alpha G_0-G_N,\qquad
 F=Q_\chi^{-1/2}\Delta^{1/2}:E\longrightarrow\PP_{q-1},\\
 R=R_N+K_\chi F.
 \end{gathered}
\end{equation}
The norm is the ordinary largest eigenvalue of the positive matrix
shown. Consequently
\(G_0^{-1/2}\Delta G_0^{-1/2}\succeq I\), so \(\Delta>0\)
and \(F\) is invertible. The orthogonal cross terms vanish and give
\begin{equation}\tag{PT29}
 J_NR=I_E,\qquad R^*H_NR=G_N+F^*Q_\chi F
                         =\alpha G_0,\qquad
 H_{\alpha G_0}=H_{G_0}.
\end{equation}
The last equality follows directly from (PT2), because a positive
scalar commutes with the specified positive square root. This is
an actual finite polynomial section realizing the trace minimum.
It retains the full relation \(\chi\), not a root-value ideal.

Its residual at every smaller polynomial source is also exact.
For \(W\subseteq\PP_{q-1}\), write
\[
 Q_W=\PP_{2q-1}/K_\chi W,\qquad
 \eta_W:\PP_{q-1}/W\hookrightarrow Q_W,\quad
 [f]\mapsto[K_\chi f],\qquad
 \pi_W:Q_W\twoheadrightarrow E.
\]
The first map is well-defined and injective because \(K_\chi\)
is injective, and the second is induced by \(J_N\).
Their image and kernel agree by \(\ker J_N=\im K_\chi\).
Therefore this is the explicit exact sequence and section comparison
\begin{equation}\tag{PT30}
 \begin{gathered}
 0\longrightarrow\PP_{q-1}/W\xrightarrow{\eta_W}Q_W
   \xrightarrow{\pi_W}E\longrightarrow0,\\
 [R(\cdot)]_W-[R_N(\cdot)]_W=\eta_W([F(\cdot)]_W),\qquad
 \ker([F]_W)=F^{-1}W,\qquad
 \rank([F]_W)=q-\dim W.
 \end{gathered}
\end{equation}
Surjectivity of \([F]_W\) follows from invertibility of \(F\).
Thus every proper finite polynomial source retains a nonzero
residual map of the computed rank. The arithmetic quotient class
is nevertheless identical, because \(\pi_W[R]_W=I_E\).

For completeness, the entire change in the multiplication boundary
is a relation as well, with its original degree recorded. Interpret
the identity below in \(\PP_{2q}\), and define
\(\widehat F:E\to\PP_q\) by \(\widehat F=SF-FA\).
Then direct substitution gives
\begin{equation}\tag{PT31}
 (SR-RA)-(SR_N-R_NA)=K_\chi^{(q)}\widehat F,
 \quad K_\chi^{(q)}:\PP_q\hookrightarrow\PP_{2q},\ f\mapsto\chi f.
\end{equation}
For a target source \(W'\subseteq\PP_q\) this difference has
precisely the class \([\widehat F]_{W'}\) under the injective
map induced by \(K_\chi^{(q)}\). Thus a full-quotient equality
has not erased the smaller-source generator residual.

\section{Backward receivers for the earlier theta source}

The original TC source has weight one and a packet of actual zeros;
the present cyclic sum packet has weight \(k\). We give the exact
receiving operations at those two types. For any positive Gram and
operator pair \((G,A_0)\) of weight \(w\), the operation is
\((G,A_0,w)\mapsto B_0=A_0-wI/2\), followed by the trace
projection proof (PT7)--(PT9) on its complete positive-primary sum.
The polynomial coordinate change to (PT1), when used, is the
original full-primary sum realization, whose explicit CRT map is
(PT3). This operation does not replace an ordered tensor space by
its cyclic image: TC21 remains on its original ordered tensor space.

In the original TC22 coordinates, write
\(W_m=-z_mb_m-b_m^*z_m^*\), where \(z_m\) is its actual
column and \(b_m\) its actual row. The receiver to our row form
is exactly
\begin{equation}\tag{PT32}
 \ell=b_m,qquad s=z_m^*,\qquad
 sG_m^{-1}s^*=u_m,quad \ell G_m^{-1}\ell^*=v_m,quad
 sG_m^{-1}\ell^*=\bar c_m.
\end{equation}
Complex conjugation changes the sign of the imaginary part but
not its square; (PT25) is therefore precisely the original scalar
certificate TC26 and lower bound TC28 in these coordinates. The
new projection formula (PT26) strengthens its description by
evaluating the excess relative to the unchanged positive-primary
exterior line. The known rank-one Gram update TC27 remains exactly
as stated; it is not assigned to the distinct sequence (PT21).
TC18 and TC21 are operator-norm infima. Equations (PT12) and
(PT18) add the positive-trace minimum and a minimizing rank-two
stratum, without changing those previously proved infima.

Here is the exact smaller-source map for the original theta
realization TC29a--d \cite{tc}, including its finite correction.
Retain
\[
 V=\{\phi\in\mathcal S(\R):\phi\text{ even},\ \phi(0)=0,
                 \ \int_\R\phi=0\},\quad D=-x\partial_x,
 \quad\Theta\phi=2\sum_{n\ge1}\phi(nx),
\]
and the original space \(\mathscr B\) of smooth functions rapidly
decreasing at both ends on \(\R_{>0}\). The map \(\Theta:V\to
\mathscr B\) is injective. Indeed on \(\Re s>1\) absolute
convergence and the substitution \(u=nx\) give
\(\mathcal M\Theta\phi(s)=2\zeta(s)\mathcal M\phi(s)\).
The Euler product is nonzero there, and uniqueness of the Mellin
transform on a vertical line is Fourier uniqueness after \(x=e^y\).
Thus \(\Theta\phi=0\) forces \(\phi=0\).
For each \(W\subseteq V\) define the actual maps
\[
 \eta_W:V/W\longrightarrow\mathscr B/\Theta W,quad
 [\phi]\longmapsto[\Theta\phi],\qquad
 \pi_W:\mathscr B/\Theta W\longrightarrow\mathscr B/\Theta V.
\]
Injectivity of \(\Theta\) proves injectivity of \(\eta_W\),
and its image is exactly \(\ker\pi_W\). For the original
reference section \(s_h\) and its finite primitive correction
\(T=\Phi_LC:E_h\to V\), put \(R=s_h+\Theta T\). Subtraction
of actual functions gives
\begin{equation}\tag{PT33}
 [R]_W-[s_h]_W=\eta_W([T]_W),\qquad
 \pi_W[R]_W=[s_h]_V,qquad
 \ker([R-s_h]_W)=T^{-1}W.
\end{equation}
The last equality uses the proved injectivity. This is the
precise map from a full-source correction to its smaller-source
class, for the original theta packet \(E_h\).

The original reference relation is
\(Ds_h-s_hA_h=\Theta(\phi_*\ell_h)\). Since termwise
differentiation gives \(D\Theta=\Theta D\), substitution proves
\begin{equation}\tag{PT34}
 DR-RA_h=\Theta K_R,\qquad
 K_R=\phi_*\ell_h+DT-TA_h:E_h\longrightarrow V.
\end{equation}
For \(DW\subseteq W'\), the induced map
\(\bar D:\mathscr B/\Theta W\to\mathscr B/\Theta W'\)
is well-defined because \(D\Theta W\subseteq\Theta W'\).
Its defect on (PT33) is exactly \(\eta_{W'}([K_R]_{W'})\),
which vanishes on a vector precisely when that original primitive
belongs to \(W'\). This proves the typed generator receiver.

For \(a>0\) let \(\mathcal R_aF(x)=F(x/a)\) and let
\(A_{a,h}=e^{(\log a)A_h}\) be the original arithmetic action.
Retain the reference primitive \(\Psi_{a,h}\) defined by
\(\mathcal R_as_h-s_hA_{a,h}=\Theta\Psi_{a,h}\).
The equality \(\mathcal R_a\Theta=\Theta\mathcal R_a\)
and substitution give
\begin{equation}\tag{PT35}
 \begin{gathered}
 c_{R,a}=\Psi_{a,h}+\mathcal R_aT-TA_{a,h},\qquad
 \mathcal R_aR-RA_{a,h}=\Theta c_{R,a},\\
 \bar{\mathcal R}_a^{W\to\mathcal R_aW}[R]_W
       -[R]_{\mathcal R_aW}A_{a,h}
       =\eta_{\mathcal R_aW}([c_{R,a}]_{\mathcal R_aW}),\\
 c_{R,ab}=\mathcal R_a c_{R,b}+c_{R,a}A_{b,h}.
 \end{gathered}
\end{equation}
For the last identity, expand the preceding two definitions and
cancel the terms \(\mathcal R_aT A_{b,h}\). The reference
primitives satisfy the same identity because the difference has
zero \(\Theta\)-image and \(\Theta\) is injective. The target
source remains \(\mathcal R_aW\), so no source invariance is
silently inserted. Equations (PT30)--(PT31) give the complete
polynomial counterpart on the unchanged sum packet, and
(PT33)--(PT35) are the actual TC proper-source receiver.

\section{Quantitative excess and the original cross metric}

The trace defect also measures the original metric interaction
between the positive and negative primary halves. In the fixed
CRT ordering let
\begin{equation}\tag{PT36}
 \begin{gathered}
 \mathcal G=(C^{-1})^*GC^{-1}
       =\begin{pmatrix}P&X\\X^*&Y\end{pmatrix},\quad
 K=P^{-1}X,\quad Z=Y-X^*P^{-1}X>0,\\
 F=\begin{pmatrix}P^{-1/2}&-KZ^{-1/2}\\0&Z^{-1/2}\end{pmatrix},
 \qquad\mathcal F=C^{-1}F,\qquad\mathcal F^*G\mathcal F=I.
 \end{gathered}
\end{equation}
Positivity of \(P\) is restriction positivity; positivity of
\(Z\) follows by testing \(\mathcal G\) on \((-Kv,v)\).
Multiplication verifies the displayed isometry. Write the original
CRT operator as \(CBC^{-1}=\diag(B_+,B_-)\), with
\(B_-=-D_-\). Its matrix in this isometry is
\begin{equation}\tag{PT37}
 \begin{gathered}
 \mathcal F^{-1}B\mathcal F
   =\begin{pmatrix}\widehat B_+&D\\0&\widehat B_-\end{pmatrix},
 \quad\widehat B_+=P^{1/2}B_+P^{-1/2},\quad
 \widehat B_-=Z^{1/2}B_-Z^{-1/2},\\
 D=P^{1/2}(KB_--B_+K)Z^{-1/2},\qquad
 \mathcal H=\mathcal F^*W(G)\mathcal F
   =\begin{pmatrix}D_+^{\rm H}&D\\D^*&D_-^{\rm H}\end{pmatrix},\\
 D_\pm^{\rm H}=\widehat B_\pm+\widehat B_\pm^*,\quad
 \Tr D_+^{\rm H}=L,\quad\Tr D_-^{\rm H}=-L.
 \end{gathered}
\end{equation}
The matrix \(G^{1/2}\mathcal F\) is unitary, so \(\mathcal H\)
and \(H_G\) have exactly the same eigenvalues. Let
\(\mathcal E=T_+(G)-L\), and let \(\norm D_1\) denote the
sum of the singular values of \(D\). Then
\begin{equation}\tag{PT38}
 \norm D_1^2\le\mathcal E(2L+\mathcal E),\qquad
 \mathcal E\ge\sqrt{L^2+\norm D_1^2}-L.
\end{equation}
To prove this, take a singular-value decomposition
\(D=U\Sigma V^*\), extending the singular frames to full unitary
bases, and put \(W=UV^*\). The matrices
\[
 J_0=\begin{pmatrix}I&0\\0&-I\end{pmatrix},\qquad
 T_0=\begin{pmatrix}0&W\\W^*&0\end{pmatrix}
\]
square to \(I\) and anticommute. Hence
\(Z_\theta=\cos\theta J_0+\sin\theta T_0\) is a Hermitian
contraction for every real \(\theta\). Diagonalizing
\(\mathcal H\), every diagonal entry of \(Z_\theta\) lies
between \(-1\) and \(1\), so
\(\Tr(Z_\theta\mathcal H)\le\Tr|\mathcal H|\).
The two sides are respectively
\(2L\cos\theta+2\norm D_1\sin\theta\) and
\(2(L+\mathcal E)\), since \(\Tr\mathcal H=0\).
Maximizing over \(\theta\) proves (PT38).

There is a complete finite inverse from this measured cross block
to the original \(X\). For positive-half block parameters
\(b_\alpha=\lambda_\alpha-c\) and negative-half parameters
\(\beta_\eta=\lambda_\eta-c\), the original operator blocks
are \(b_\alpha I+J_e\) and \(\beta_\eta I+J_e\).
Set \(R=P^{-1/2}DZ^{1/2}\). Then
\begin{equation}\tag{PT39}
 K_{\alpha\eta}
 =-\sum_{u,v=0}^{e-1}
 \frac{(-1)^u(u+v)!}
      {u!v!(b_\alpha-\beta_\eta)^{u+v+1}}
 J_e^uR_{\alpha\eta}J_e^v,qquad X=PK.
\end{equation}
Indeed \(R=KB_--B_+K\) by (PT37). The absolutely convergent
integral \(-\int_0^\infty e^{-tB_+}Re^{tB_-}\,dt\) solves
this equation by integration of its derivative. Each real
exponent is at most \(-2\delta t\), multiplied by a finite
polynomial. Expansion of its two full nilpotent exponentials
gives (PT39), with sign \((-1)^u\) because the right exponential
is \(e^{+tB_-}\). Uniqueness follows from the homogeneous
equation: \(e^{-tB_+}Ke^{tB_-}=K\) then holds for all \(t\),
while its left side tends to zero. Thus the integral and the
finite sum invert the original cross map exactly.

For an explicit bound in the same full primary coordinates set
\begin{equation}\tag{PT40}
 \begin{gathered}
 \kappa=\max_{\alpha,\eta}
       \sum_{u,v=0}^{e-1}
       \frac{(u+v)!}{u!v!|b_\alpha-\beta_\eta|^{u+v+1}},\\
 \norm{P^{-1}X}_F^2
 \le\kappa^2\norm{P^{-1/2}}^2\norm{Z^{1/2}}^2
                  \mathcal E(2L+\mathcal E).
 \end{gathered}
\end{equation}
Each power of \(J_e\) has norm at most one. Applying the triangle
inequality to every block in (PT39), then summing squared block
Frobenius norms, gives \(\norm K_F\le\kappa\norm R_F\).
The multiplication bound for \(R=P^{-1/2}DZ^{1/2}\), followed
by \(\norm D_F\le\norm D_1\) and (PT38), proves (PT40).
All original metric factors are displayed; no uniform bound on
their condition numbers has been inserted.
For the canonical metric (PT21) one substitutes exactly
\(G=G_N\) and \(\mathcal E=\epsilon_N-L\). Thus the same
scalar which occurs in (PT25)--(PT26) controls the actual
positive/negative primary metric cross block through (PT36)--(PT40).
In particular the original scalar certificate itself strengthens to
\begin{equation}\tag{PT41}
 u_Nv_N-(\Im z_N)^2\ \ge\ L^2+\norm{D_N}_1^2,
 \qquad
 D_N=P_N^{1/2}(K_NB_--B_+K_N)Z_N^{-1/2},
\end{equation}
where \(P_N,X_N,Y_N,K_N,Z_N\) are obtained by substituting
the unchanged \(G_N\) in (PT36). Indeed
\(\mathcal E(2L+\mathcal E)=(\epsilon_N-L)
(\epsilon_N+L)=\epsilon_N^2-L^2\), so (PT38) and (PT25)
give precisely (PT41). This is an estimate on the actual attained
metric; every factor in its cross block has its original coordinate
map and its finite inverse in (PT36)--(PT39).

\section{Scope of the finite result}

The positive-trace minimum (PT12), every minimizing metric
(PT14)--(PT16), the rank-two attainer (PT18), the exterior growth
(PT19)--(PT20), the canonical scalar and defect (PT25)--(PT26),
the finite source and proper-source residuals (PT28)--(PT35),
and the complete cross-metric receiver (PT36)--(PT41)
are established by the displayed proofs. The fixed attained
metrics remain the matrices in (PT21). The variational minimum
does not evaluate their separate projected kernel and boundary
allocations; their exact trace gap is (PT26). No value for an
unevaluated projection is assigned in this result.

\begin{thebibliography}{9}
\bibitem{incoming}
User-supplied web continuation, \emph{Exact minimum of the full-packet
positive weight trace}, 19 September 2026,
\texttt{04\_FULL\_METRIC\_POSITIVE\_TRACE\_MINIMUM.md}, M1--M7.
Exact local source hashes and paths accompany this note.
\bibitem{tc}
Split-Zero research archive, \emph{Exact source-metric optimization and
aggregate arithmetic control}, \texttt{tau\_chain.tex}, TC1--TC29d,
in the 13 September 2026 recursive integration repository.
\bibitem{gmf}
Split-Zero mathematical continuation, \emph{Growing marked flags in
the original arithmetic quotient metrics}, 19 September 2026,
\texttt{GROWING\_MARKED\_FLAGS.tex}, GMF1--3.
\bibitem{fan}
Ky Fan, ``On a theorem of Weyl concerning eigenvalues of linear
transformations. I,'' \emph{Proceedings of the National Academy of
Sciences of the United States of America} \textbf{35} (1949), 652--655.
\href{https://doi.org/10.1073/pnas.35.11.652}{doi:10.1073/pnas.35.11.652}.
The finite trace variational argument used here is proved in full
in (PT7)--(PT9).
\bibitem{lyapunov}
A.~M.~Lyapunov, \emph{The general problem of the stability of motion},
1892; English translation in \emph{International Journal of Control}
\textbf{55} (1992), no.~3. Publisher record for the translated work:
\href{https://doi.org/10.1080/00207179208934253}{doi:10.1080/00207179208934253}.
The finite matrix integral, positivity criterion and complete
nilpotent evaluation used here are proved in (PT13)--(PT16).
\end{thebibliography}
\end{document}
